% !TeX program = xelatex
\documentclass[margin = 5mm, tikz]{standalone}
\usepackage{xfrac}
\usepackage[MnSymbol]{mathspec}
\usepackage{ifthen}
\usepackage{esvect}
\setmainfont[Mapping=tex-text,Ligatures={NoRequired,NoCommon,NoContextual}]{Comic Sans MS}
\setmathfont(Digits,Latin,Greek)[Numbers={Lining,Proportional}]{Comic Sans MS}
\usetikzlibrary{shapes.geometric,decorations,decorations.pathmorphing}
\usetikzlibrary{intersections}
\usetikzlibrary{patterns}
\usetikzlibrary{calc}
\usetikzlibrary{shapes.misc}
\usetikzlibrary{spy}

\pgfdeclarelayer{bg}
\pgfdeclarelayer{bbg}
\pgfsetlayers{bbg, bg, main}

\definecolor{c1}{HTML}{ac0000}
\definecolor{c2}{HTML}{00a0d5}

\newcommand{\abs}[1]{{\left|{#1}\right|}}

\tikzset{
	point/.style n args = {1}{circle, draw = black, inner sep = #1pt, fill = black},
	empty point/.style = {circle, draw = black, inner sep = 2pt, fill = white},
	cross/.style = {cross out, draw, minimum size = 2 * (#1 - \pgflinewidth), inner sep = 0pt, outer sep = 0pt},
	xkcd/.style n args = {1}{decorate, decoration = {random steps, pre length = #1mm, post length = #1mm, segment length = 0.6mm, amplitude = 0.2pt}},
	xkcd/.default = {4},
	point/.default = {2}
}

\begin{document}
	\begin{tikzpicture}
		\pgfmathsetmacro{\rw}{0.2}
		\draw[xkcd, thick] (-3,0) -- (7,0);
		\foreach \i in {-2,0,...,6} {
			\pgfmathtruncatemacro{\labeli}{\i/2}
			\ifthenelse{\i=2}{
				\node[point = {1.8}, c1, label = {below}:{\small$\labeli$}] (a\i) at (\i,0) {};
			}{
				\node[point = {1.8}, label = {below}:{\small$\labeli$}] (a\i) at (\i,0) {};
			}
		}
		\node[point = {1.8}, c2, label = {above, c2}:{\small$\sqrt{2}$}] (sqrt2) at ({sqrt(8)},0) {};
		\node[point = {1.8}, c1, label = {below, c1}:{\small$\sfrac{3}{2}$}] (sqrt2) at (3,0) {};
		\fill[xkcd, fill opacity = 0.1, c1] (-3,-\rw) rectangle (2,\rw);
		\fill[xkcd, fill opacity = 0.1, c1] (3,-\rw) rectangle (7,\rw);
		\fill[xkcd, fill opacity = 0.1, c2] (2,-\rw) rectangle (3,\rw);
	\end{tikzpicture}
\end{document}